%% ma: L12-B17-0007
%% lop: 12
%% chuong: 5
%% bai: 17
%% dang: 12.17.2
%% loai: DS
%% muc: [TH, TH, TH, VD]
%% dap_an: "ĐĐĐS"
%% nguon: "(NBV)-ĐỀ SỐ 8-MỖI NGÀY 1 ĐỀ THI 2025.pdf (D08, MỖI NGÀY 1 ĐỀ - Đề số 8), Phần II Câu 2"
%% kien_thuc: "Bài 17 (khoảng cách giữa hai điểm, mặt cầu), Bài 8 (tọa độ trong không gian)"
%% trang_thai: chinh-thuc
%% ghi_chu: "Đáp án gốc Đ Đ Đ S đúng (sympy: M(0;1;1) duy nhất). Cùng mô hình với L12-B17-0002 (dạng 12.17.2: bốn khoảng cách), dạng Đúng/Sai. Gốc chỉ có ảnh vệ tinh trang trí, đã bỏ."
\begin{ex}
Trong không gian với hệ toạ độ $Oxyz$, cho bốn vệ tinh $A(0;4;5)$, $B(0;5;4)$, $C(1;3;3)$, $D(1;-1;3)$. Điểm $M(a;b;c)$ trong không gian, biết khoảng cách từ các vệ tinh đến điểm $M$ lần lượt là $AM=5$, $BM=5$, $CM=3$, $DM=3$.
\choiceTF{\True $a^2+(b-4)^2+(c-5)^2=a^2+(b-5)^2+(c-4)^2=25$}{\True $(a-1)^2+(b-3)^2+(c-3)^2=(a-1)^2+(b+1)^2+(c-3)^2=9$}{\True $b=c$}{$M(1;1;1)$}
\loigiai{a) Đúng: $AM^2=BM^2=25$ cho $a^2+(b-4)^2+(c-5)^2=a^2+(b-5)^2+(c-4)^2=25$.

b) Đúng: $CM^2=DM^2=9$ cho $(a-1)^2+(b-3)^2+(c-3)^2=(a-1)^2+(b+1)^2+(c-3)^2=9$.

c) Đúng: từ $(b-4)^2+(c-5)^2=(b-5)^2+(c-4)^2$ suy ra $-8b-10c=-10b-8c$, tức $b=c$.

d) Sai: từ $(b-3)^2=(b+1)^2$ được $b=1$, nên $c=1$; thay vào $a^2+(b-4)^2+(c-5)^2=25$ được $a^2=0$, $a=0$ (thỏa $(a-1)^2+4+4=9$). Vậy $M(0;1;1)$, không phải $M(1;1;1)$.}
\end{ex}
